% page 480 du fac-simile (image p-479 du scan)
% bloc 162.6 x 246.3 mm ; marge gauche 8.0 mm
\pgc{-12.700mm}{0.000mm}{
\l{\cel{3}472}
\l{}
\l{pulular. (\sou{netrans.}) Kreskar (nombrale) abundege e rapide.}
\l{\cel{3}- EFIS.}
\l{}
\l{\sou{pulvero.}\cel{2}I. Substanco solida, medikamenta, reduktita a par-}
\l{\cel{3}tikuli tre tenua. - II. (\sou{milit-arto}). Mixuro de sulfo, sal-}
\l{\cel{3}petro e karbono, reduktita a partikuli tenua, qua, kande}
\l{\cel{3}acendata, esas extreme exploziva - ed uzesas en la pafarmi}
\l{\cel{3}por lansar projektili. - DEFIS.}
\l{}
\l{\sou{pumao.}\cel{2}(\sou{zool.}) Mamifero karnavida, ek la familio \textquotedbl{}felidi\textquotedbl{},}
\l{\cel{3}qua vivas sur la kontinento Amerika. - L.\cel{1}\sou{felix concolor}.}
\l{\cel{3}- DEFIRS.}
\l{}
\l{\sou{pumico}. (\sou{geol.)}\cel{2}Petro volkanala, poroza,lejera, uzata por}
\l{\cel{3}frotar, polisar, e c. - EIS.}
\l{}
\l{\sou{pumpar}. (\sou{trans.}) Aspirar o pulsar (liquido, gaso) per aparato}
\l{\cel{3}qua, ordinare, konsistas ye cilindro en qua su movas, rekta-}
\l{\cel{3}linee ed alterne, pistono. - DEFIS.}
\l{}
\l{\sou{puncar}. (\sou{trans.}) Surpozar (quaze perforo-probe) marko konkava}
\l{\cel{3}produktita per stalo-bloko kun, ye un ek lua extremaji pinta,}
\l{\cel{3}graburo reliefa, e sur la altra extremajo di qua onu frapas}
\l{\cel{3}per martelo. - II. Perforar (bilieto kartona o papera) per}
\l{\cel{3}pinco por igar lu nevalida. - DEFIS.}
\l{}
\l{\sou{puncho}. Drinkajo, mixuro ek aquo, od infuzuro de teo, kun}
\l{\cel{3}brandio o rumo, citrono e sukro. - DEFIRS.}
\l{}
\l{\sou{puncionar}. (\sou{trans.}) (\sou{kirurgio}).Pikar kirurgiale por igar es-}
\l{\cel{3}kapar liquido qua amaseskis aden kavajo di la korpo (kava-}
\l{\cel{3}jo naturala od acidentala). - DEFIS.}
\l{}
\l{\sou{punisar}. (\sou{trans.}) (\sou{pro, per)}.\cel{2}Aplikar (ad ulu) sufrigilo pro}
\l{\cel{3}lua kulpo. - EFIS.}
\l{}
\l{\sou{punsar}. (\sou{trans.)}\cel{2}Kopiar desegnuro per pulvero quan onu pas-}
\l{\cel{3}igas tra la trui di papero perforita segun la linei di la}
\l{\cel{3}desegnuro. - EF.}
\l{}
\l{\sou{punto}. I. (\sou{gram}.)Signo ortografiala super le \textquotedbl{}i\textquotedbl{}. - II.}
\l{\cel{3}+Puntuo-signo:\cel{2}. ; ? ! ...\cel{2}III. (\sou{muziko}).Signo qua, apud}
\l{\cel{3}noto, augmentas la duro-valoro di olca ; qua indikas halt-}
\l{\cel{3}instanto en la mezuro. - IV. (\sou{geom.}) Elemento qua genitas la}
\l{\cel{3}lineo, konsiderata per abstrakto, kom ulo qua esas nek longa,}
\l{\cel{3}nek larja, nek profunda. - V. La subdividuro minima di la}
\l{\cel{3}parti di la temo quan onu traktas. - VI. (\sou{ludi : biliardo},}
\l{\cel{3}\sou{karti, loto, e}\cel{1}c.)\cel{2}Marko quan onu facas sur paperfolio od}
\l{\cel{3}ardezo-plako, pos singla gano-stroko. -VII.(\sou{tipogr.)}\cel{1}Mezuro}
\l{\cel{3}qua determinas la grandeso di la tipi. - DEFIRS.}
\l{}
\l{pupo.\cel{2}(navig.)\cel{2}Dopa parto di navo, di batelo. - EFIS.}
\l{}
\l{pupeo. Mikra figuro homala, ek\cel{1}\sou{kar}tono, vaxo, ligno, uzata}
\l{\cel{3}da la infantini kom ludilo. - DEF.}
}
